\scanpage{253}
\margin{Grothendieck}
\mspage{252}
n$^\circ$ 293

\uline{LE THEOREME PRINCIPAL DE Z A R I S K I}

I. - LEMMES PRELIMINAIRES

\uline{Lemme} 1 - \uline{Soient} $r$ \uline{un domaine d'intégrité et} $o$ \uline{un sous-anneau de} $r$. \uline{On suppose qu'il existe des éléments} $t_1,\ldots,t_n$ \uline{de} $r$ $(n>0)$ \uline{qui sont algébriquement indépendants par rapport à} $o$ \uline{et tels que} $r$ \uline{soit entier sur} $o[t_1,\ldots,t_n]$. \uline{Soient} $q$ \uline{un idéal premier maximal de} $r$ \uline{et} $p=q\cap o$. \uline{Il y a alors un idéal premier de} $r$ \uline{contenu dans} $q$, \uline{distinct de} $q$ \uline{et contenant} $p$.

Nous nous ramènerons d'abord au cas où $o$ est intégralement clos. Soit $o^*$ la clôture intégrale de $o$, et soit $r^*=r[o^*]$; les éléments $t_1,\ldots,t_n$ sont algébriquement indépendants sur $o^*$ et $r^*$ est entier sur $o^*[t_1,\ldots,t_n]$. Soit $q^*$ un idéal premier de $r^*$ tel que $q^*\cap r=q$; comme $r^*$ est entier sur $r$, $q^*$ est un idéal premier maximal de $r^*$. Soit $p^*=q^*\cap o^*$. Supposons qu'il y ait un idéal premier $n^*$ de $r^*$ contenu dans $q^*$, distinct de $q^*$ et contenant $p^*$. Soit $n=n^*\cap r$; $n$ est un idéal premier de $r$; comme $q^*\cap r=q$ on a $n\subset q$; il résulte du fait que $n^*\ne q^*$ et de ce que $r^*$ est entier sur $r$ que l'on a $n\ne q$. Comme $n^*$ contient $p^*=q^*\cap o^*$, il contient aussi $q\cap o=p$, d'où $p\subset n$.

Supposons donc maintenant que $o$ soit intégralement clos. On sait qu'il en est de même de $o[t_1,\ldots,t_n]$. L'idéal $po[t_1,\ldots,t_n]$ engendré par $p$ dans $o[t_1,\ldots,t_n]$ est premier. Soit en effet $R$
